\documentclass[11pt,a4paper]{article}

\usepackage[utf8]{inputenc}
\usepackage[T1]{fontenc}
\usepackage{amsmath,amssymb}
\usepackage{booktabs}
\usepackage{graphicx}
\usepackage[margin=2.5cm]{geometry}
\usepackage[colorlinks=true,linkcolor=blue,urlcolor=blue]{hyperref}

\graphicspath{{../analysis/results/cnn/}}

\newcommand{\code}[1]{\nolinkurl{#1}}
\newcommand{\csitl}{CsI(Tl)}
\setlength{\emergencystretch}{2em}

\title{How to reproduce the CryspLight CNN spatial-resolution study\\
\large Direct three-dimensional regression and reconstructability selection
in \csitl{} and BGO}
\author{J.J.~G\'omez-Cadenas}
\date{July 22, 2026}

\begin{document}
\maketitle

\begin{abstract}
This note is an operational guide to the CNN spatial-resolution analysis already
completed in CryspLight. It records the dependency order, commands,
configurations, selections and outputs needed to reproduce two results for both
\csitl{} and BGO: direct regression of the first-interaction position
$(x,y,z)$, and the same response after retaining the 80\% of events judged most
reconstructable by a learned classifier. It does not prescribe regenerating
existing data or models: every stage begins with the artifact that can be
reused. The classifier used in the published result is the
\emph{reconstructability classifier}, not the exploratory single-site versus
multi-site topology classifier.
\end{abstract}

\tableofcontents

%=====================================================================
\section{What is being reproduced}

The input to every estimator is the time-integrated $8\times8$ SiPM
photoelectron map plus the logarithm of the total number of photoelectrons. The
target is the true position of the first gamma interaction in the crystal. The
two crystals and realistic photopeak selections are

\begin{center}
\begin{tabular}{lccc}
\toprule
crystal & dimensions [mm$^3$] & energy FWHM at 511 keV & window \\
\midrule
\csitl{} & $48\times48\times37.2$ & 6\%  & $511\pm2\sigma$ keV \\
BGO      & $48\times48\times22.4$ & 10\% & $511\pm2\sigma$ keV \\
\bottomrule
\end{tabular}
\end{center}

Two cases are measured on held-out test events:

\begin{enumerate}
  \item \textbf{Direct regression.} CryspNet maps each selected event directly
  to $(x,y,z)$. The realistic energy window contains direct photoelectric and
  multi-interaction Compton events.
  \item \textbf{Regression after classification at 80\% acceptance.} A second
  CryspNet assigns an event-quality score. The positive training label is
  ``the direct regressor misses the first interaction by at least 5 mm in
  three dimensions.'' Keeping the lowest-scoring 80\% produces the selected
  response. The direct regressor is not retrained for this published result;
  the classifier is an event-selection layer in front of its test predictions.
\end{enumerate}

The directory also contains classifiers named \code{cnn\_*\_cls}. They predict
truth topology, with \code{n\_int >= 2} labelled multi-site. Those classifiers
were an exploratory comparison and are \emph{not} used for the 80\% resolution
reported here. The production classifiers are \code{cnn\_*\_rcls}, where
\code{rcls} means reconstructability classification.

%=====================================================================
\section{Environment and artifact inventory}

Run commands from the repository root. The supported setup is Apple Silicon
macOS with Julia 1.11 or newer and Python 3.11 or newer. Instantiate the locked
environments only if they are not already available:

\begin{verbatim}
julia --project=. -e 'using Pkg; Pkg.instantiate()'
python3 -m venv .venv
.venv/bin/python -m pip install -r analysis/requirements.txt
\end{verbatim}

Before running an expensive stage, check for its terminal artifact:

\begin{center}
\begin{tabular}{p{0.27\textwidth}p{0.63\textwidth}}
\toprule
stage & reusable terminal artifact \\
\midrule
500k event simulation &
\code{output/gammas\_\{csitl,bgo\}\_500k/events.h5} \\
direct regression &
\code{analysis/results/cnn/cnn\_\{csitl,bgo\}\_win/cnn\_best.pt} and
\code{test\_events.h5} \\
two-site comparison &
\code{analysis/results/cnn/cnn\_\{csitl,bgo\}\_2site/test\_events.h5} \\
reconstructability classifier &
\code{analysis/results/cnn/cnn\_\{csitl,bgo\}\_rcls/cnn\_best.pt} and
\code{test\_events.h5} \\
80\% response analysis &
\code{analysis/results/cnn/selection\_fits/selection\_fits.json} \\
PET response recipe &
\code{analysis/results/cnn/selection\_fits/pet\_resolution\_recipe.json} \\
\bottomrule
\end{tabular}
\end{center}

If an artifact exists and its upstream configuration has not changed, start at
the next stage. In particular, calculating tables or figures from the committed
\code{test\_events.h5} files takes seconds and requires neither gamma simulation
nor CNN training.

%=====================================================================
\section{Stage 0: event datasets}

The simulation configurations are
\code{runs/gammas\_csitl\_500k.toml} and
\code{runs/gammas\_bgo\_500k.toml}. Each shoots 500,000 incident 511 keV
gammas and writes the light map, detected photoelectrons and interaction truth.
The common photodetection efficiency is 0.40. Seeds, materials, dimensions,
surface model and binning are fixed in the TOML files.

\textbf{Skip this stage} when both expected \code{events.h5} files exist. To
recreate only missing datasets, run the corresponding command:

\begin{verbatim}
julia --project=. --threads=auto scripts/shoot_gammas.jl \
    runs/gammas_csitl_500k.toml

julia --project=. --threads=auto scripts/shoot_gammas.jl \
    runs/gammas_bgo_500k.toml
\end{verbatim}

The outputs are \code{output/gammas\_csitl\_500k/events.h5} and
\code{output/gammas\_bgo\_500k/events.h5}. This is the expensive optical
transport stage.

%=====================================================================
\section{Stage 1: direct coordinate regression}
\label{sec:direct}

The production configurations are
\code{analysis/cnn/runs/cnn\_csitl\_win.toml} and
\code{analysis/cnn/runs/cnn\_bgo\_win.toml}. The suffix \code{win} identifies
the realistic energy-window mixture. Do not substitute \code{cnn\_csitl.toml}
or \code{cnn\_bgo.toml}: those configurations select idealized, contained
direct-photoelectric events and answer a different question.

To retrain only when the saved checkpoints or test dataframes are absent:

\begin{verbatim}
.venv/bin/python analysis/cnn/train.py \
    analysis/cnn/runs/cnn_csitl_win.toml

.venv/bin/python analysis/cnn/train.py \
    analysis/cnn/runs/cnn_bgo_win.toml
\end{verbatim}

The driver performs the seeded window smearing, a seeded 70/15/15
train/validation/test split, exact D4 augmentation of the training split, CNN
and MLP training, checkpoint selection on validation RMSE, and final test
evaluation. Its principal outputs are

\begin{itemize}
  \item \code{metrics.json}: coordinate RMSE, coordinate $p_{68}$ and the
  fraction with three-dimensional error above 5 mm, globally and by event type;
  \item \code{test\_events.h5}: event row key, truth, energy, topology and all
  estimator predictions; this is the reusable input to later selections;
  \item \code{cnn\_best.pt} and \code{history.json}: checkpoint and training
  history;
  \item \code{residuals.png}, \code{rmse\_vs\_z.png},
  \code{distance\_by\_type.png} and \code{convergence.png}: diagnostic figures.
\end{itemize}

The outputs live below \code{analysis/results/cnn/cnn\_csitl\_win/} and
\code{analysis/results/cnn/cnn\_bgo\_win/}. The already obtained headline
test results are

\begin{center}
\begin{tabular}{lccc}
\toprule
crystal & RMSE $(x/y/z)$ [mm] & $p_{68}(|\Delta x|/|\Delta y|/|\Delta z|)$ [mm]
& $P(|\Delta\mathbf r|>5\ \mathrm{mm})$ \\
\midrule
\csitl{} & $3.19/3.24/4.16$ & $2.02/2.03/2.08$ & 26.8\% \\
BGO      & $1.62/1.62/2.46$ & $1.08/1.08/1.26$ & 9.7\% \\
\bottomrule
\end{tabular}
\end{center}

For the unselected two-Gaussian core--tail description, use the saved
test dataframes; no training is involved:

\begin{verbatim}
.venv/bin/python analysis/cnn/fit_residuals.py \
    cnn_csitl_win cnn_bgo_win
\end{verbatim}

This writes \code{residual\_fit.json} and \code{residual\_fit.png} inside each
run directory. The core widths are approximately 1.1--1.2 mm in \csitl{} and
0.8--1.0 mm in BGO.

%=====================================================================
\section{Stage 2: prerequisite comparison products}

The classifier driver evaluates its score against the predicted separation of
a six-output, two-site network. Consequently it expects the two-site
\code{test\_events.h5} files even though those networks are not inputs to the
classifier training itself. If the files are absent, create them with

\begin{verbatim}
.venv/bin/python analysis/cnn/train.py \
    analysis/cnn/runs/cnn_csitl_2site.toml

.venv/bin/python analysis/cnn/train.py \
    analysis/cnn/runs/cnn_bgo_2site.toml
\end{verbatim}

If \code{analysis/results/cnn/cnn\_*\_2site/test\_events.h5} already exists,
skip this stage. The two-site regression is a diagnostic comparator; the
selected spatial response still uses the three-output predictions from
Section~\ref{sec:direct}.

%=====================================================================
\section{Stage 3: reconstructability classifier}
\label{sec:classifier}

The production classifier configurations are
\begin{itemize}
  \item \code{analysis/cnn/runs/cnn\_csitl\_rcls.toml};
  \item \code{analysis/cnn/runs/cnn\_bgo\_rcls.toml}.
\end{itemize}
For every event, the saved direct regressor supplies a predicted
first-interaction position. The classifier label is
\begin{equation}
 y_{\rm bad}=\mathbb{1}\!\left(
 \|\widehat{\mathbf r}_1-\mathbf r_1\|\geq5\ {\rm mm}\right).
\end{equation}
The classifier is trained to assign high scores to these badly reconstructed
events. A selection therefore accepts events below a score threshold.

To recreate missing classifier artifacts:

\begin{verbatim}
.venv/bin/python analysis/cnn/train_classifier.py \
    analysis/cnn/runs/cnn_csitl_rcls.toml

.venv/bin/python analysis/cnn/train_classifier.py \
    analysis/cnn/runs/cnn_bgo_rcls.toml
\end{verbatim}

These commands require the \code{cnn\_*\_win} checkpoints and the
\code{cnn\_*\_win} and \code{cnn\_*\_2site} test dataframes. They write
\code{cnn\_best.pt}, \code{test\_events.h5}, \code{metrics.json} and
\code{classification.png} below \code{analysis/results/cnn/cnn\_*\_rcls/}.
The test dataframe carries the original event row and classifier score, allowing
an exact join to the regression dataframe. The obtained classifier AUC is
0.796 for \csitl{} and 0.793 for BGO.

The topology configurations \code{cnn\_*\_cls.toml} must not be used here.
Their AUCs, 0.719 and 0.640, quantify single/multi-site discrimination, but
topology is not identical to reconstruction quality: close multi-site events
are often reconstructed accurately and should not be rejected.

%=====================================================================
\section{Stage 4: the 80\% selected response}
\label{sec:selection}

The response calculation is an analysis of existing held-out predictions:

\begin{verbatim}
.venv/bin/python analysis/cnn/selection_fits.py
\end{verbatim}

For each crystal the script joins \code{cnn\_*\_win/test\_events.h5} with
\code{cnn\_*\_rcls/test\_events.h5} by the original event row. It considers
100\%, 80\%, 60\% and 50\% global acceptance, retaining the events with the
lowest reconstructability scores. For each coordinate it fits
\begin{equation}
 p(\Delta q)=f_{\rm core}\,
 \mathcal N(\mu,\sigma_1)+(1-f_{\rm core})\,
 \mathcal N(\mu,\sigma_2), \qquad q\in\{x,y,z\}.
\end{equation}
The no-selection fit determines $\mu$, the core width $\sigma_1$ and tail width
$\sigma_2$. In the comparison used for the resolution prescription these
shapes are frozen and only the two amplitudes are refitted after selection.
This tests whether classification removes the broad component instead of
artificially redefining or narrowing the core.

The outputs are

\begin{itemize}
  \item \code{analysis/results/cnn/selection\_fits/selection\_fits.json}, the
  complete machine-readable result;
  \item \code{selection\_fits.txt}, a compact table of free-fit and fixed-shape
  amplitudes;
  \item \code{selection\_fits.png}, all residual distributions and fits.
\end{itemize}

At 80\% acceptance the measured fraction of events with a 3D error above 5 mm
falls from 26.8\% to approximately 17.1\% in \csitl{}, and from 9.7\% to 5.6\%
in BGO. The fixed-shape depth tail-to-core amplitude ratio changes from 0.486 to
0.114 in \csitl{} and from 0.218 to 0.043 in BGO. The core widths remain those
of the unselected response.

%=====================================================================
\section{Stage 5: compact resolution recipe}

To turn the fitted response into parameters for a system-level PET simulation,
run

\begin{verbatim}
.venv/bin/python analysis/cnn/pet_recipe.py
\end{verbatim}

The script reads \code{selection\_fits.json} and writes
\code{pet\_resolution\_recipe.json} in the same directory. For each crystal,
coordinate and selection scenario it provides $\sigma_1$, $\sigma_2$ and
$f_{\rm core}$. The scenarios included are no classifier selection, 80\%
acceptance and 60\% acceptance. Acceptance is a per-photon efficiency; requiring
both photons in a coincidence to pass gives the square of that efficiency.

%=====================================================================
\section{Minimal rerun recipes}

\subsection*{Only regenerate the published tables and figures}

When all committed/saved test dataframes are present, run only

\begin{verbatim}
.venv/bin/python analysis/cnn/fit_residuals.py \
    cnn_csitl_win cnn_bgo_win
.venv/bin/python analysis/cnn/selection_fits.py
.venv/bin/python analysis/cnn/pet_recipe.py
\end{verbatim}

\subsection*{Retrain classifiers but reuse regression and simulation}

Provided the \code{cnn\_*\_win} and \code{cnn\_*\_2site} products exist:

\begin{verbatim}
.venv/bin/python analysis/cnn/train_classifier.py \
    analysis/cnn/runs/cnn_csitl_rcls.toml
.venv/bin/python analysis/cnn/train_classifier.py \
    analysis/cnn/runs/cnn_bgo_rcls.toml
.venv/bin/python analysis/cnn/selection_fits.py
.venv/bin/python analysis/cnn/pet_recipe.py
\end{verbatim}

\subsection*{Rebuild the full analysis from existing events.h5 files}

Run Stages 1 through 5 in order. Stage 0 remains unnecessary. The seeds in all
TOML configurations make the energy smearing and dataset partitions
reproducible; the event-row join guarantees that regression and classification
are compared on the same held-out events.

%=====================================================================
\section{Interpretation checklist}

\begin{itemize}
  \item Quote both the narrow core and broad tail; a single RMS obscures the
  response structure of the realistic Compton mixture.
  \item Treat 80\% as event acceptance, not an 80\% classification accuracy or
  an 80\% geometric tolerance.
  \item Use \code{rcls}, not \code{cls}, for the published selection.
  \item Do not describe the selected result as a newly trained 80\% regressor.
  It is the direct regressor evaluated after an event-quality selection.
  \item Keep the ±5 mm residual definition and energy-window parameters fixed
  when comparing crystals or reruns.
  \item For provenance, retain the TOML configuration, checkpoint,
  \code{metrics.json} and \code{test\_events.h5} together under the run tag.
\end{itemize}

The full scientific motivation, architecture and interpretation are in
\code{latex/crysplight\_cnn.tex}. This how-to is deliberately narrower: it is
the executable map from existing artifacts to the direct and 80\%-selected
spatial-resolution results.

\end{document}
